\section{Conclusion and Future Work}

This study introduces a three-stage machine learning model to assess the effects of fossil fuel price volatility on the WACC of energy transition projects. The findings emphasize that achieving the 1.5$^\circ$C target requires region-specific climate finance solutions. Policymakers and investors need to tailor their strategies to address the unique challenges posed by fossil fuel price volatility, since this can promote renewable energy investments that remain financially viable and resilient amid global market shifts. The results also reinforce the SDG framing introduced at the outset: de-risking clean-energy finance supports affordable and clean energy (SDG 7), climate action (SDG 13), open and reproducible innovation infrastructure (SDG 9), and cross-regional research partnerships (SDG 17).

While this study provides some insights into the financial dynamics of the energy transition, further research is also needed. Initially, some datasets and assumptions rely on secondary sources. Future work could enhance the robustness of these findings by incorporating firm-level financial data and longer regional panels. Additionally, future work could transition from quantifying predictive sensitivity to formal causal inference, using techniques such as fixed effects, instrumental variables, or Structural Equation Models to handle policy variables, financial mechanisms, and technological learning rates. Finally, the framework itself can be enhanced by benchmarking against advanced architectures such as Transformers or state-space models and propagating uncertainty across all three stages, including the GPR predictive variance, for uncertainty quantification and stress-testing.
